\section{The Translation Ladder}
\label{sec:ladder}

The phrase ``translate what a dog says'' smuggles in an assumption: that a
sentence-shaped meaning waits in the signal, the way a French sentence waits
behind its English translation. The biological record gives reason to doubt
this. We therefore replace the single verb with a ladder of five distinct
capabilities, each with its own input--output contract, its own notion of
ground truth, and its own distance from the present
(Figure~\ref{fig:ladder}).

\begin{figure}[t]
\centering
\resizebox{\textwidth}{!}{% Figure 1 — The Translation Ladder (hero figure).
% Compiled standalone via fig1_standalone.tex; also \input by sections/03_ladder.tex.
\begin{tikzpicture}[
  font=\sffamily,
  >={Latex[length=2.6mm,width=2mm]},
  block/.style={rounded corners=3pt,draw=ink!70,line width=0.7pt,
    minimum height=0.9cm,align=left,inner sep=5pt,text width=2.8cm,
    font=\sffamily\small},
  num/.style={font=\sffamily\bfseries\large},
  note/.style={font=\sffamily\scriptsize\itshape,text=ink!70},
  pill/.style={rounded corners=6pt,fill=cream,draw=rosedeep,line width=0.7pt,
    inner sep=4pt,font=\sffamily\scriptsize\bfseries,text=rosedeep,align=center},
]
  \node[block,fill=sky!45] (r1) at (0,0)
    {\num{1}\;\;\textbf{Signal}\\[1pt]\scriptsize detection \& segmentation};
  \node[block,fill=sage!55] (r2) at (3.2,1.5)
    {\num{2}\;\;\textbf{Affect}\\[1pt]\scriptsize emotion / arousal class};
  \node[block,fill=gold!55] (r3) at (6.4,3.0)
    {\num{3}\;\;\textbf{Intent}\\[1pt]\scriptsize goal + context inference};
  \node[block,fill=paw!60] (r4) at (9.6,4.5)
    {\num{4}\;\;\textbf{Reference}\\[1pt]\scriptsize symbol $\rightarrow$ object};
  \node[block,fill=rose!55] (r5) at (12.8,6.0)
    {\num{5}\;\;\textbf{Dialogue}\\[1pt]\scriptsize open-ended exchange};

  \node[note,rotate=25] at (3.4,4.9) {increasing capability $\cdot$ increasing difficulty};

  \draw[dashed,rosedeep,line width=1pt] (-0.55,2.25) -- (14.4,2.25);
  \node[pill] at (2.7,2.25)
    {THE SEMANTIC GAP\\[1pt]\mdseries no generative grammar to translate};

  \node[rounded corners=5pt,fill=sagedeep,text=white,inner sep=4pt,
    font=\sffamily\scriptsize\bfseries,align=center] (here) at (4.7,0.35)
    {YOU ARE HERE\\[1pt]\mdseries consumer products + most ML};
  \draw[->,sagedeep,line width=1pt] (here.north) -- (r2.south);
  \begin{scope}[shift={(2.78,1.74)},scale=0.5]
    \fill[sagedeep] (0,0) ellipse (0.16 and 0.13);
    \fill[sagedeep] (-0.16,0.18) circle (0.06);
    \fill[sagedeep] (0,0.23) circle (0.06);
    \fill[sagedeep] (0.16,0.18) circle (0.06);
  \end{scope}

  \node[note,align=center,anchor=west] at (8.1,2.55)
    {soundboards reach here\\(with human symbols)};
  \draw[->,ink!45,line width=0.7pt] (8.05,2.6) -- (r3.south east);

  \draw[decorate,decoration={brace,amplitude=5pt,mirror},sagedeep,line width=0.8pt]
    ($(r1.south west)+(0,-0.7)$) -- ($(r2.south east)+(0,-0.7)$)
    node[midway,below=7pt,note,text=sagedeep,align=center]
      {data-rich, falsifiable\\[1pt]\bfseries achievable now};
  \draw[decorate,decoration={brace,amplitude=5pt},rosedeep,line width=0.8pt]
    ($(r4.north west)+(0,0.55)$) -- ($(r5.north east)+(0,0.55)$)
    node[midway,above=7pt,note,text=rosedeep,align=center]
      {data-poor, partly ill-posed\\[1pt]\bfseries aspirational};
\end{tikzpicture}}
\caption{\textbf{The translation ladder.} Current consumer products and most
machine-learning work occupy \rung{Rung 2} (affect classification);
controlled soundboard studies reach \rung{Rung 3} (referential intent) but
only by supplying a human-designed symbol set. The dashed boundary marks the
\emph{semantic gap}: above it, the dog's natural signal lacks the recursive,
generative syntax that sentence-level translation presupposes, so the target
becomes possibly ill-posed rather than merely hard.}
\label{fig:ladder}
\end{figure}

\textbf{Rung 1, signal detection and segmentation}, asks only whether there
is a vocalization, where it begins and ends, and who produced it. This is a
classical detection problem with objective ground truth, and it is
essentially solved: bioacoustic detectors segment and attribute calls
reliably~\cite{ghani2023perch,hagiwara2023aves}, and individual
identification from bark acoustics is above
chance~\cite{molnar2006discriminate}. Nothing about Rung 1 requires
understanding.

\textbf{Rung 2, affect and arousal classification}, maps a signal to a small
set of affective categories or to continuous arousal/valence dimensions.
This is where nearly all commercial ``translators'' live and where the
peer-reviewed ML literature is strongest, because the labels are
human-perceived emotions or eliciting contexts that annotators can agree
on~\cite{pongracz2006acoustic,icassp2018dogbark,sci2022explainable}. Rung 2
is falsifiable and genuinely useful for welfare monitoring. Its limit is
that an emotion label is not a sentence; rendering ``fearful'' as ``I am
scared of the vacuum'' adds narrative the signal never contained.

\textbf{Rung 3, intent and context inference}, asks what the dog is trying
to do, given the surrounding context, and requires multimodal fusion and a
model of the dog's goals. Soundboard studies provide the cleanest evidence,
because the button press makes the intent explicit and
checkable~\cite{bastos2024plos,bastos2024scirep}, but the symbol set is
human-designed and trained. Rung 3 on natural signals is far less developed
and much harder to validate.

\textbf{Rung 4, referential semantics}, would map a spontaneous canine
signal to a specific referent without human-supplied symbols. Receptive
studies show dogs can understand such
mappings~\cite{kaminski2004rico,pilley2011chaser}, but a productive, open
referential vocabulary in natural dog communication has not been
demonstrated. This is the first rung where the translation framing strains:
there may be no stable symbol-to-referent dictionary to recover.

\textbf{Rung 5, open-ended dialogue}, is the folkloric dream: free,
generative, two-way conversation. It presupposes a recursive generative
grammar on the canine side, and the evidence of graded affective signals,
linear associative combinations, and no demonstrated syntax gives no
support for that presupposition. Rung 5 is therefore not merely distant; in
its strong form it is possibly ill-posed. A weaker engineered version,
structured human-to-dog input via soundboards with the dog responding
through trained symbols, is conceivable and already nascent.

The ladder converts a vague question into answerable ones: which rung, with
what ground truth, at what validated accuracy? It also explains the rhetoric
of the market, where vendors demonstrate Rung 2 while marketing Rung 5. The
rest of the survey uses it as a ruler.
